\textbf{Những công việc đã thực hiện trong đồ án}

% {\color{red} CHỈ GHI LÀ KẾT LUẬN, KHÔNG GHI CHƯƠNG 4. KẾT LUẬN NHÉ}

Đồ án đã thiết kế được phương pháp ước lượng vị trí điểm trọng tâm, xây dựng đội hình và lên chiến thuật đẩy vật dựa vàod điều khiên lực.  Đồ án đã hoàn thành các mục tiêu đề ra, đạt được các kết quả cụ thể sau đây:
\begin{itemize}
    \item Về lý thuyết:
    
    Đồ án đã trình bày rõ các lý thuyết về cơ học vật rắn, grasping, các phương pháp mô hình hóa và giải các bài toán tối ưu về cấu hình và lực, từ đó xây dựng một khung phân tích hoàn chỉnh cho bài toán tổng hợp đội hình robot nhằm sinh wrench mong muốn tại khối tâm vật thể.
    
    \item Về thực nghiệm:
    
    % Thực hiện xem xét đánh giá các thuật toán lập đường đi bao phủ và chỉ ra ưu điểm và hạn chế của của thuật toán lập đường đi đề xuất. 
    Đồ án đã thực hiện mô phỏng thuật toán với các thí nghiệm khác nhau, trong đó cấu trúc vật thể gồm đa giác lồi và không lồi bằng ngôn ngữ lập trình Python, từ đó kết quả cho thấy khả năng khám phá và thao tác trên bất kỳ hình dạng nào của vật thể dựa vào phương pháp đề xuất.
    
    Như vậy thuật toán sẽ hoạt động hiệu quả ở với hình dạng vật thể và số lượng robot tham gia bất kỳ, dễ dàng tùy biến để hoạt động với các đầu vào khác nhau. Tuy nhiên thuật toán vẫn còn các nhược điểm cần được tối ưu như quy trình điều khiển còn đơn giản, phương pháp ước lượng trọng tâm còn nhạy với nhiễu.   
\end{itemize}
    





